The universal approximation theorem guarantees arbitrarily accurate approximation of continuous functions on compact domains with suitable activations and sufficient width. How should an optimizer find the weights, and which architecture should it use? We compare weight-space programming (WSP) with gradient optimization on ten algorithmic tasks. The fixed-network study covers 29 nonconvex task--architecture pairs and three random starts, using 3--50 coefficients in rectified linear, residual and recurrent networks. Astra, a pretrained language model, returns numerical weights in one reply or revises them using training and development feedback. Both examples-only WSP procedures match every extrapolation reference in the same 30 of 87 settings, including addition and running sums. Full-batch gradient descent using backpropagation achieves this in zero settings, as does descent followed by limited-memory Broyden--Fletcher--Goldfarb--Shanno (L-BFGS) refinement. Summing each task's mean raw extrapolation root mean squared error (RMSE) ranks the two iterative WSP conditions first among 14 evaluated pipelines. The examples-only sum falls from $1.06\times10^{12}$ to $2.03\times10^9$, driven mainly by improved Fibonacci predictions. The median scaled extrapolation error instead increases, and stored square-network coefficients show how central validation can favor a worse outer continuation. Thirty further requests let Astra choose the graph and coefficients from input/output examples, with extrapolation as the primary objective. Twenty-one proposals are exact on every outer input. Square uses a single product node with zero stored coefficients, while sine attains RMSE $7.79\times10^{-16}$. Two self-power proposals are invalid. We report coefficient storage, live tensor memory and arithmetic costs under the declared operation library, and reproduce prompts, returned graphs and public construction explanations.

